\documentclass[10pt,a4paper]{amsart}
\usepackage[margin=19mm]{geometry}
\usepackage{amsmath,amssymb,mathtools}
\usepackage[scr=rsfs,cal=euler]{mathalfa}
\usepackage{xfrac}
\usepackage[unicode,hidelinks,bookmarks=false]{hyperref}
\newcommand{\ie}{i.e.\ }
\newcommand{\cf}{cf.\ }
\newcommand{\resp}{respectively\ }
\newcommand{\Subsection}{\subsection}
\providecommand{\nobreakdash}{\nobreak\hbox{-}}
\setlength{\emergencystretch}{2em}
\multlinegap0pt
\usepackage{amssymb}
\newtheorem{proposition}{Proposition}
\title{Exotic embedded surfaces and involutions from Real Seiberg--Witten theory}
\author{David Baraglia}
\date{Source: arXiv:2504.00281v2 (CC BY 4.0)}
\begin{document}
\maketitle

\noindent\emph{Source and licence.} Exotic embedded surfaces and involutions from Real Seiberg--Witten theory. Version arXiv:2504.00281v2 (CC BY 4.0), distributed by arXiv under the Creative Commons Attribution 4.0 International licence, http://creativecommons.org/licenses/by/4.0/, as recorded in the arXiv licence metadata for this exact version and shown on the abstract page https://arxiv.org/abs/2504.00281v2. Full licence notice: \texttt{licenses/arxiv-2504.00281-CC-BY-4.0.txt}.\par\smallskip

\noindent\textit{Editorial scope.} Counted unit: the article's proposition proposition (source0.tex lines 773--775). The article's complete original proof of it is delivered with it (source0.tex lines 776--810, one \texttt{proof} environment of 35 lines). No mathematics is written, repaired or re-derived: the statement and the proof are the article's own lines, in the article's own notation, and no retained line is substituted or re-flowed.  No further source-local context is needed: every cross-reference of every retained line resolves inside the two delivered spans.  Nothing was omitted from any retained span, so the package carries no omission record.  No retained line contains a citation, so the wrapper carries no bibliography and no bibliography entry.  No retained line contains a figure environment, an image-inclusion command or drawing code, so no image is delivered and the package declares no asset dependency.  Editorial formatting only, in addition: an AMS \texttt{amsart} preamble that restates the article's own declarations and macros used by the retained lines, namely the environments \texttt{proposition} on separate counters, together with the standard packages those lines need.  The retained environments print in this document as Proposition 1 in order of appearance, whereas the article prints the same items with the numbering of its own full document.  The complete native source of the article as distributed in the arXiv e-print for this exact version is bundled unchanged as \texttt{sources/175738/source0.tex}.
\begin{proposition}
Suppose $d$ is even and $w_1(D_R) = 0$. Then $deg_R(X,\mathfrak{s})$ does not depend on the choice of splitting.
\end{proposition}

\begin{proof}
Since any two splittings $s_1,s_2 : H^1(X ; \mathbb{Z})^{-\sigma} \to \mathcal{H}_R$ differ by a homomorphism $H^1(X ; \mathbb{Z})^{-\sigma} \to \mathbb{Z}_2$, they agree on the subgroup $2 H^1(X ; \mathbb{Z})^{-\sigma} \subseteq H^1(X ; \mathbb{Z})^{-\sigma}$. This means that there is a canonical lift $s : 2 H^1(X ; \mathbb{Z})^{-\sigma} \to \mathcal{H}_R$. Now we imitate the usual construction of the Real Bauer--Furuta map, except that instead of quotienting by $s(H^1(X ; \mathbb{Z})^{-\sigma})$, we quotient by $s(2H^1(X;\mathbb{Z})^{-\sigma})$. The result is a $\Gamma$-equivariant map 
\[
\widetilde{f} : S^{\widetilde{V} , \widetilde{U}} \to S^{\widetilde{V}' , \widetilde{U}'}
\]
over $\widetilde{Jac}_R(X) = H^1(X;\mathbb{R})^{-\sigma}/2H^1(X;\mathbb{Z})^{-\sigma}$, where $\Gamma = \mathcal{H}_R/s(H^1(X;\mathbb{Z})^{-\sigma}$. Observe that $\Gamma$ is a $\mathbb{Z}_2$ central extension of $\mathbb{Z}_2^{\beta}$, where $\beta = b_1(X)^{-\sigma}$ and $\widetilde{Jac}_R(X) \to Jac_R(X)$ is a $\mathbb{Z}_2^\beta$-covering space. Splittings $s' : \mathbb{Z}_2 \to \Gamma$ are in bijection with splittings $s' : H^1(X ; \mathbb{Z})^{-\sigma} \to \mathcal{H}_R$ (there is a natural map from splittings $H^1(X;\mathbb{Z})^{-\sigma} \to \mathcal{H}_R$ to splittings $\mathbb{Z}_2^\beta \to \Gamma$ and this map is easily seen to be a bijection). Given a splitting $s : \mathbb{Z}_2^\beta \to \Gamma$, the Bauer--Furuta map corresponding to $s$ is obtained by quotienting $\widetilde{f}$ by $s(\mathbb{Z}_2^\beta)$. Set $\widetilde{D} = \widetilde{V} - \widetilde{V}'$. Then the virtual bundle $D_R$ corresponding to a chosen splitting $s$ is given by $D_R = \widetilde{D}/s(\mathbb{Z}_2^\beta)$. Since $\widetilde{f}$ is $\Gamma$-equivariant, it has a $\Gamma$-equivariant degree
\[
deg_{\Gamma}( \widetilde{f} ) \in H^{b_+(X)^{-\sigma}-d}_\Gamma( \widetilde{Jac}_R(X) ; \mathbb{Z}).
\]
Pulling back under a splitting $s : \mathbb{Z}_2^\beta \to \Gamma$, we obtain
\[
s^*( deg_{\Gamma}(\widetilde{f})) \in H^{b_+(X)^{-\sigma}-d}_{\mathbb{Z}_2^\beta}( \widetilde{Jac}_R(X) ; \mathbb{Z}) \cong H^{b_+(X)^{-\sigma}-d}(Jac_R(X) ; \mathbb{Z}).
\]
Clearly $s^*(deg_{\Gamma}(\widetilde{f}))$ coincides with $deg(f)$, where $f$ is the Bauer--Furuta invariant corresponding to $s$. So it remains to show that $s^*( deg_{\Gamma}(\widetilde{f}))$ does not depend on the choice of splitting.

Fix a splitting $s_0$, which we use to identify $\Gamma$ with $\mathbb{Z}_2 \times \mathbb{Z}_2^\beta$. Under this splitting
\[
H^*_{\Gamma}(\widetilde{Jac}_R(X) ; \mathbb{Z}) \cong H^*_{\mathbb{Z}_2}( Jac_R(X) ; \mathbb{Z}).
\]
We have an isomorphism
\[
H^*_{\mathbb{Z}_2}( Jac_R(X) ; \mathbb{Z}) \cong H^*( Jac_R(X) ; \mathbb{Z}) \otimes H^*_{\mathbb{Z}_2}(pt ; \mathbb{Z}).
\]
Recall that $H^*_{\mathbb{Z}_2}(pt ; \mathbb{Z}) \cong \mathbb{Z}[x]/(2x)$, where $deg(x) = 2$.

Now let $s : \mathbb{Z}_2^\beta \to \Gamma$ be any splitting. Then the pullback
\[
s^* : H^*_{\mathbb{Z}_2}( Jac_R(X) ; \mathbb{Z}) \to H^*( Jac_R(X) ; \mathbb{Z})
\]
is the map
\[
H^*( Jac_R(X) ; \mathbb{Z}) \otimes \mathbb{Z}[x]/(2x) \to H^*(Jac_R(X) ; \mathbb{Z})
\]
which sends $x$ to some class $a = s^*(x) \in H^*(Jac_R(X) ; \mathbb{Z})$. Since $2x=0$, we must also have $2a=0$. But $H^*(Jac_R(X) ; \mathbb{Z})$ has no $2$-torsion and hence $s^*( deg_{\Gamma}(\widetilde{f}))$ does not depend on the choice of $s$.
\end{proof}

\end{document}
